
\nopagebreak

La región está limitada por la curva \(y=5\sqrt{x}\), el eje \(x\) y las rectas
\(x=1\) y \(x=4\).  

Dividimos el intervalo \([1,4]\) en \(n\) partes iguales:
\[
\Delta x = \frac{4-1}{n} = \frac{3}{n}.
\]

Tomando puntos derechos:
\[
x_i = 1 + \frac{3i}{n}.
\]

El área se expresa como límite de sumas de Riemann:
\[
A = \lim_{n\to\infty} \sum_{i=1}^{n} 5\sqrt{x_i} \, \Delta x.
\]

Sustituyendo \(x_i\):
\[
A = \lim_{n\to\infty} \sum_{i=1}^{n} 5 \sqrt{1 + \frac{3i}{n}} \cdot \frac{3}{n}.
\]

Simplificando:
\[
A = \lim_{n\to\infty} \sum_{i=1}^{n} \frac{15}{n} \sqrt{1 + \frac{3i}{n}}.
\]

Este es el **límite de sumas de Riemann** que representa el área exacta.  

Opcionalmente, evaluando la integral:
\[
A = \int_{1}^{4} 5\sqrt{x}\,dx
= 5 \int_{1}^{4} x^{1/2}\,dx
= 5 \left[ \frac{2}{3} x^{3/2} \right]_1^4
= \frac{10}{3} \left( 8 - 1 \right)
= \frac{70}{3}.
\]

Finalmente:
\[
\boxed{A = \frac{70}{3}}
\]
